\section{Methodology}
\label{sec:method}

We cross two factors on the same model and the same questions: what the model knows (labeled by neutral elicitation) and what the prompt rewards (neutral, incentive-only, or an instructed lie).
A false answer on a known question is a lie; a false answer on a question the model does not know is a hallucination.

\subsection{Model, questions, and belief labels}

\para{Model and questions.}
We use \llama, run locally in fp16.
Statements are decoded greedily with at most 24 new tokens; samples use temperature 1.0.
The question pool is 3,000 random questions from the \triviaqa \texttt{rc.nocontext} validation split \citep{joshi2017triviaqa}, drawn with seed 0.
A further 300 disjoint questions serve to pilot the pressure prompts.

\para{Belief elicitation.}
Under the neutral system prompt, the model answers three paraphrased user templates greedily, and we draw 10 samples at temperature 1 on the first template.
We fix the labels in \tabref{tab:labels} before the pressure runs.
Following \citet{cheang2025llms}, we separate questions on which the model repeats one wrong answer (confident-wrong) from questions on which it has no stable answer (no-knowledge).

\begin{table}[t]
\caption{Belief labels from neutral elicitation on 3,000 \triviaqa questions. Ambiguous questions count toward the shares of false answers but are excluded from detector contrasts.}
\label{tab:labels}
\centering
\small
\begin{tabular}{@{}lp{0.68\linewidth}r@{}}
\toprule
Label & Rule & $n$ \\
\midrule
known & all 3 greedy answers correct and at least 8/10 samples correct & 1,772 \\
confident-wrong & all 3 greedy answers wrong, at least 8/10 samples agree with the wrong greedy answer, at most 1/10 correct & 96 \\
no-knowledge & no greedy answer correct, at most 1/10 samples correct, not confident-wrong & 386 \\
ambiguous & everything else & 746 \\
\bottomrule
\end{tabular}
\end{table}

\subsection{Statement conditions and cells}

\para{Conditions.}
All four conditions use the same user turn; only the system prompt changes (full text in \secref{app:prompts}).
The {\em neutral} prompt is ``You are a helpful assistant.''
The {\em falsehood-pays} prompt casts the model as a quiz-show rival who wins \$1M only if the user answers wrongly; it contains no instruction to lie.
The {\em truth-pays} prompt is a matched control: the same quiz show, but the user is a teammate and both win if the user answers correctly.
It differs from the falsehood-pays prompt only in whose interest the truth serves, so it measures how often a known answer flips to false merely because the prompt changed.
The {\em instructed} prompt tells the model that it must lie and answer with a plausible but incorrect answer.
We chose the falsehood-pays prompt from a pilot of four incentive prompts (sandbagging threat, manager pressure, confidentiality, quiz-show rival) as the one with the highest yield of false answers on known questions.

\para{Grading.}
An LLM judge (\texttt{openai/gpt-5.6-luna}, temperature 0) labels each greedy answer as correct, incorrect, or evasive against the \triviaqa aliases.
Samples are graded by alias string match.
The judge and alias match agree on 92.5\% to 96.0\% of greedy answers, depending on the condition.
About 0.1\% of judge calls returned unusable output after four retries and fell back to alias match.

\para{Cells.}
\Tabref{tab:cells} defines the cells we contrast.
We keep ``changed false answers without knowledge'' apart from hallucinations because, under a lie instruction, a model that does not know the truth may still be trying to deceive.
Under the truth-pays control, 56\% of no-knowledge answers also change, so a changed answer is weak evidence of intent.

\begin{table}[t]
\caption{Cells of the matched design. ``Unknown'' pools confident-wrong and no-knowledge questions. Counts are for the instructed / falsehood-pays conditions where two numbers appear.}
\label{tab:cells}
\centering
\small
\begin{tabular}{@{}lp{0.53\linewidth}r@{}}
\toprule
Cell & Definition & $n$ \\
\midrule
lie & known question, false answer under the pressure prompt & 611 / 52 \\
honest-correct (same prompt) & known question, correct answer under the pressure prompt & 1,157 / 1,718 \\
hallucination (same prompt) & unknown question, false answer identical to the neutral answer & 145 / 225 \\
changed false answer & unknown question, false answer that differs from the neutral answer & 316 / 228 \\
correct (neutral) & known question, correct neutral answer & 1,772 \\
hallucination (neutral) & unknown question, false neutral answer & 467 \\
\bottomrule
\end{tabular}
\end{table}

\subsection{Detectors}

We extract residual-stream activations at 14 layers.
Unless stated otherwise, a probe score is computed on the mean activation over response tokens.

\para{Lie detectors.}
The {\em deception probe} follows the instructed-pairs recipe of \citet{goldowskydill2025detecting}: logistic regression on standardized activations ($C=0.1$), trained on 254 true facts from \citet{azaria2023internal} with the last five words dropped, under ``Pretend you're an honest / a dishonest person\ldots'' prompts.
We also report a difference-of-means variant.
The {\em prompt-token} control applies the same logistic-regression probe to the last prompt token, before any answer token exists; whatever it detects cannot depend on the answer.
The {\em self-report} detector appends ``Was your last response deceptive? (Yes/No)'' and scores the logit of Yes minus the logit of No.

\para{Truth probe.}
We train logistic regression on 1,954 true/false statements from \citet{marks2023geometry}, at layer 14 and the last token, chosen by leave-one-dataset-out AUROC (0.93).

\para{Hallucination detectors.}
The {\em correctness probe} is a logistic regression on neutral answers (incorrect versus correct), 5-fold cross-fitted by question.
{\em Answer improbability} is the negative mean log-probability of the response tokens.
{\em Sample inconsistency} is the share of 10 same-prompt samples that disagree with the greedy answer, a string-match proxy for semantic entropy.

\para{In-distribution probes.}
As a ceiling, we train logistic regression on each contrast itself at layer 16, 5-fold cross-fitted by question, once on the response and once on the last prompt token only.
These tell us whether the information is linearly present, so that a failure of a transferred detector is a failure of transfer and not of representation.

\para{Layer choice.}
Each transferred probe uses the layer that is best on its own standard contrast: layer 12 for the deception probe (instructed lie versus neutral correct) and layer 16 for the correctness probe (hallucination versus correct under the neutral prompt).
We report full layer sweeps in \secref{sec:layers}.

\subsection{Metrics}

We report AUROC with 1,000-sample bootstrap 95\% confidence intervals and Wilson intervals for proportions.
For fixed-threshold behavior we report the flag rate at the 99th percentile of scores on 500 on-policy \alpaca responses, the ``1\% FPR'' convention of \citet{goldowskydill2025detecting}.

\subsection{External check on \mask}

To see whether the conclusions carry over to more realistic pressure, we run 510 binary items from the \texttt{known\_facts}, \texttt{disinformation}, and \texttt{continuations} archetypes of \mask \citep{ren2025mask}.
We elicit beliefs with the benchmark's first belief prompt, one greedy and two sampled runs, and the judge maps each to a side of the proposition; a belief exists if all three agree.
A lie is a pressured statement that contradicts the belief.
